% =============================================================================
%  Section 2.4 — Các công trình liên quan
% =============================================================================

\section{Các công trình liên quan}

Để giải quyết bài toán biểu diễn và truy vấn tri thức, nhiều hướng tiếp cận đã được đề xuất trong những năm gần đây. Dựa trên đặc thù kỹ thuật và phương pháp tích hợp tri thức, các nghiên cứu đi trước có thể được phân loại thành ba nhóm chính. Các nhóm này đều đạt được các cột mốc đáng kể trong việc hỗ trợ giảm sự mơ hồ trong tri thức của LLMs tuy nhiên mỗi nhóm đều bộc lộ những hạn chế cụ thể khi áp dụng vào miền tri thức nghệ thuật có cấu trúc phức tạp như nghệ thuật Chèo.

Nhóm nghiên cứu đầu tiên tập trung vào các hệ thống hỏi đáp dựa trên đồ thị tri thức (KGQA) và ứng dụng đồ thị tri thức (KG) trong biểu diễn di sản văn hóa. Các công trình của Fan et al. \cite{fan2023cichmkg}, Wan et al. \cite{wan2024wumkg} và Huang et al. \cite{huang2023using} đã khẳng định khả năng ưu việt của KG trong việc kết nối các nguồn dữ liệu văn hóa phân mảnh. Tuy nhiên, hạn chế cốt lõi của các dự án này là mới dừng lại ở khâu tổ chức và lưu trữ dữ liệu, thiếu sự tích hợp với các mô hình sinh ngôn ngữ tự nhiên để tạo thành một hệ thống hỏi đáp hoàn chỉnh. Mặt khác, đối với các hệ thống KGQA truyền thống, các nghiên cứu dựa trên semantic parsing như Sun et al. (2020)\cite{sun2020sparqa} và Yin et al. (2016) \cite{yih2015semantic} đã chỉ ra rằng phương pháp này phụ thuộc mạnh vào việc phân tích cú pháp và các mẫu truy vấn được định nghĩa sẵn. Sự phụ thuộc này khiến hệ thống trở nên kém linh hoạt và gặp khó khăn khi xử lý các truy vấn phức tạp hoặc yêu cầu suy luận đa bước.

Nhóm nghiên cứu thứ hai liên quan đến kiến trúc Retrieval-Augmented Generation (RAG) do Lewis et al. \cite{lewis2020retrieval} khởi xướng. Phương pháp này giải quyết được bài toán cập nhật tri thức cho LLMs mà không cần huấn luyện lại. Mặc dù vậy, RAG truyền thống tồn tại nhược điểm lớn về "truy xuất phẳng". Bằng cách chia nhỏ tài liệu thành các đoạn văn bản độc lập và so khớp dựa trên độ tương đồng vector (embedding), cơ chế này vô tình cắt đứt cấu trúc quan hệ nội tại giữa các thực thể. Do đó, đối với các truy vấn yêu cầu tổng hợp thông tin từ nhiều nguồn hoặc theo dấu các liên kết logic phức tạp, RAG thường thất bại trong việc cung cấp một ngữ cảnh đủ sâu và liền mạch.

Nhóm nghiên cứu thứ ba tập trung vào kiến trúc Graph-based RAG (GraphRAG), ra đời nhằm khắc phục hạn chế của RAG bằng cách tích hợp trực tiếp cấu trúc đồ thị vào quá trình truy xuất. Các công trình như KG-guided RAG \cite{zhu2025knowledge} hay GRAG \cite{hu2025grag} đã sử dụng KG để định hướng việc mở rộng ngữ cảnh một cách hiệu quả. Tiến xa hơn, các hệ thống SG-RAG \cite{saleh2025sgrag} và KRAGEN \cite{matsumoto2024kragen} tập trung truy xuất các tiểu đồ thị (subgraph), trong khi phương pháp của Gong et al. \cite{gong2025beyond} khai thác sâu bộ ba quan hệ (triplet-driven). Hiện tại, các hệ thống GraphRAG hiện hành vẫn vấp phải một rào cản kỹ thuật lớn: chúng bị giới hạn ở cơ chế truy xuất đơn độ hạt (single-granularity). Việc cố định một mức độ truy xuất làm hệ thống thiếu tính linh hoạt khi phải xử lý đa dạng các loại câu hỏi, từ định danh thực thể cục bộ đến phân tích toàn cục. Thêm vào đó, các kiến trúc này chủ yếu được phát triển cho các miền tri thức phổ quát hoặc y sinh. Khi áp dụng vào nghệ thuật sân khấu dân gian—nơi hệ thống nhân vật và lớp diễn có sự biến thiên liên tục—chúng bộc lộ sự thiếu tương thích trầm trọng do thiếu một Ontology chuyên biệt.

Từ những khoảng trống nghiên cứu và hạn chế kỹ thuật của các công trình đi trước, sự cần thiết của một hệ thống tối ưu hóa riêng cho nghệ thuật Chèo trở nên vô cùng rõ ràng. Khóa luận này đóng góp vào hướng nghiên cứu GraphRAG bằng cách đề xuất một kiến trúc truy xuất đa độ hạt (multi-granularity retrieval) linh hoạt—chuyển đổi mượt mà từ mức độ node, triplet, path đến subgraph. Cách tiếp cận này không chỉ vượt qua giới hạn của RAG truyền thống và GraphRAG đơn độ hạt, mà còn bảo tồn trọn vẹn cấu trúc tự sự của nghệ thuật Chèo, đảm bảo tính trung thực và khả năng truy vết trong quá trình LLM tạo sinh câu trả lời.
